% P15.30: Durrett Exercises 4.3.4--4.3.10, with their actual proof routes.
\chapter{Event Products and Elementary Density Constructions}
\label{sec:p15-density-and-product-exercises}

Infinite products connect event occurrence, survival probabilities and
equivalence of sequence laws. Finite partitions also give an elementary
density construction on a countably generated sigma-field. This chapter
supplies all the source arguments in Exercises~4.3.4--4.3.10
\cite[pp.~233--234]{DurrettPTEAuthor2019}, including the
partition-based Radon--Nikodym proof with no dependence on the general
existence theorem for conditional expectation.

\section{Event products and survival}

\begin{hbtheorem}{thm:p15-event-product-criterion}{proved=annex:p15.event-product-criterion}{The Borel--Cantelli proof of the product criterion}
For real numbers \(0\leq p_m<1\),
\[
                 \prod_{m=1}^\infty(1-p_m)>0
                     \quad\Longleftrightarrow\quad
                 \sum_{m=1}^\infty p_m<\infty.
\]
The product means the decreasing limit of its finite products.
The annex supplies the source's independent-event and Borel--Cantelli
route, separately from the logarithmic proof of
\cref{thm:p15-product-affinity-series}.
\end{hbtheorem}

\begin{hbtheorem}{thm:p15-survival-product}{proved=annex:p15.survival-product}{Conditional hazards without independent events}
Let \(A_1,A_2,\ldots\) be events, and put
\(D_n=\bigcap_{j=1}^nA_j^c\).
If \(P(D_n)=0\) for any finite \(n\), then
\(P(\bigcap_{j\geq1}A_j^c)=0\).
Otherwise define the scalar conditional hazards
\[
                     h_n=P(A_n\mid D_{n-1})\quad(n\geq2).
\]
They lie in \([0,1)\), and
\[
 P(D_n)=P(A_1^c)\prod_{j=2}^n(1-h_j),\qquad
 P\!\left(\bigcap_{j\geq1}A_j^c\right)=0
             \ \Longleftrightarrow\ \sum_{j=2}^\infty h_j=\infty.
\]
No independence of the events is required. Conditioning on a
zero-probability survival event is not part of this statement.
\end{hbtheorem}

\section{Finite partitions and a density theorem}

\begin{hbdefinition}{def:p15-partition-averages}{Elementary conditional averages}
Let \(\nu\) be a probability measure and
\(\mathcal P_n\) finite measurable partitions with each partition
refining its predecessor. Put \(\mathcal F_n=\sigma(\mathcal P_n)\).
For an integrable \(W\), its elementary conditional average at a
cell \(I\in\mathcal P_n\) with \(\nu(I)>0\) is
\[
                         \frac{\int_I W\,d\nu}{\nu(I)}.
\]
Assign zero on the \(\nu\)-null cells.
This finite sum of cell constants is a measurable integrable function
and has the same integral as \(W\) on every event of \(\mathcal F_n\).
This definition uses only finite arithmetic and integration.
\end{hbdefinition}

\begin{hbtheorem}{thm:p15-partition-ratio-supermartingale}{proved=annex:p15.partition-ratio-supermartingale}{Refinement ratios with and without local absolute continuity}
Let \(\mu\) be a finite nonnegative measure, with the probability
\(\nu\) and refining partitions just defined. Set
\[
 X_n|_I=
 \begin{cases}\mu(I)/\nu(I),&\nu(I)>0,\\0,&\nu(I)=0.\end{cases}
\]
Then \(X_n\geq0\), \(E_\nu X_n\leq\mu(\Omega)\), and its elementary
conditional averages satisfy
\(E_\nu(X_{n+1}\mid\mathcal F_n)\leq X_n\).
If \(\mu|_{\mathcal F_n}\ll\nu|_{\mathcal F_n}\) at every time,
equality holds and \(X_n\) is a martingale with constant integral
\(\mu(\Omega)\).

Without local absolute continuity, strict loss occurs, for example,
when \(\mu=\delta_0\), \(\nu\) is Lebesgue probability on \([0,1]\),
\(\mathcal P_0=\{\Omega\}\), and all later partitions are
\(\{\{0\},(0,1]\}\): \(X_0=1\) and \(X_n=0\), \(\nu\)-almost surely,
for all \(n\geq1\).
\end{hbtheorem}

\begin{hblemma}{thm:p15-elementary-partition-convergence}{proved=annex:p15.elementary-partition-convergence}{Bounded convergence from finite averaging alone}
Suppose each \(Y_n\) is constant on cells of \(\mathcal P_n\),
\(0\leq Y_n\leq M<\infty\) almost surely, and elementary conditional
averaging gives \(E_\nu(Y_{n+1}\mid\mathcal F_n)=Y_n\).
Then \(Y_n\) has a measurable limit \(Y\in[0,M]\) almost surely
and converges in \(L^1(\nu)\). Its proof uses finite predictable
trading sums and the upcrossing argument, without assuming a
Radon--Nikodym representation for any measure.
\end{hblemma}

\begin{hbtheorem}{thm:p15-countably-generated-radon-nikodym}{proved=annex:p15.countably-generated-radon-nikodym}{The source's elementary Radon--Nikodym construction}
Let \(\mathcal F=\sigma(E_1,E_2,\ldots)\), let \(\mu,\nu\) be
sigma-finite nonnegative measures on \((\Omega,\mathcal F)\), and
assume \(\mu\ll\nu\). Then there is a nonnegative measurable \(f\)
with
\[
                            \mu(A)=\int_A f\,d\nu
                                   \quad(A\in\mathcal F).
\]
It is unique \(\nu\)-almost everywhere and finite
\(\nu\)-almost everywhere. The annex constructs it from the finite
partitions generated by \(E_1,\ldots,E_n\), using only their
elementary averages and the preceding bounded-convergence lemma.
It supplies both the finite/probability reduction and the complete
sigma-finite assembly requested in source Exercise~4.3.7.
\end{hbtheorem}

\section{All Bernoulli product cases}

\begin{hbdefinition}{def:p15-bernoulli-sequence-laws}{Two sequences of success probabilities}
For \(\alpha_m,\beta_m\in[0,1]\), let
\(\mu=\bigotimes_{m\geq1}\operatorname{Bernoulli}(\alpha_m)\) and
\(\nu=\bigotimes_{m\geq1}\operatorname{Bernoulli}(\beta_m)\)
on the binary sequence space, where \(\xi_m\) is its \(m\)th coordinate.
Define
\[
 a_m=\sqrt{\alpha_m\beta_m}
                 +\sqrt{(1-\alpha_m)(1-\beta_m)}\in[0,1].
\]
In this chapter \(a_m\) is an affinity and the Greek parameters
are scalar success probabilities. These laws also have their
canonical realization on real sequence space, concentrated on binary
sequences, so the preceding real-coordinate product theorem applies.
\end{hbdefinition}

\begin{hbtheorem}{thm:p15-bernoulli-product-complete}{proved=annex:p15.bernoulli-product-complete}{The exact criterion, including endpoint support}
The full absolute-continuity criterion is
\[
 \mu\ll\nu
 \quad\Longleftrightarrow\quad
 \left\{
 \begin{array}{l}
  \beta_m=0\Longrightarrow\alpha_m=0,\quad
  \beta_m=1\Longrightarrow\alpha_m=1
                   \quad\text{for every }m,\\
  \prod_{m\geq1}a_m>0 .
 \end{array}\right.
\]
The implications on the first line are the local support condition.
Under that condition every \(a_m>0\), and the product condition
is equivalent to \(\sum_m(1-a_m)<\infty\).

If \(0<\epsilon\leq1/2\) and every
\(\alpha_m,\beta_m\in[\epsilon,1-\epsilon]\), then
\[
 \frac{(\alpha_m-\beta_m)^2}{4(1-\epsilon)}
       \leq1-a_m
       \leq\frac{(\alpha_m-\beta_m)^2}{4\epsilon}.
\]
In this regime \(\mu\) and \(\nu\) are equivalent exactly when
\(\sum_m(\alpha_m-\beta_m)^2<\infty\), and are singular otherwise.
The full endpoint criterion is an explicit extension of the source's
positive-coordinate setting.
\end{hbtheorem}

\begin{hbtheorem}{thm:p15-bernoulli-sparse-singularity}{proved=annex:p15.bernoulli-sparse-singularity}{A tail event separating two laws with small squared differences}
If \(\sum_m\alpha_m<\infty\) and \(\sum_m\beta_m=\infty\), then
\(\mu\perp\nu\). The event of finitely many success coordinates
has \(\mu\)-probability one and \(\nu\)-probability zero.

For example \(\alpha_m=(m+1)^{-2}\) and
\(\beta_m=(m+1)^{-1}\), \(m\geq1\), are all strictly between zero
and one and satisfy \(\sum_m(\alpha_m-\beta_m)^2<\infty\),
but their product laws are singular.
The uniform bound away from zero and one in the preceding quadratic
criterion cannot be dropped.
\end{hbtheorem}

\begin{hbtheorem}{thm:p15-bernoulli-l1-comparison}{proved=annex:p15.bernoulli-l1-comparison}{Summable absolute differences suffice without uniform interior bounds}
If \(0<\alpha_m,\beta_m<1\) for every \(m\) and
\(\sum_m|\alpha_m-\beta_m|<\infty\), then the two product laws
are mutually absolutely continuous.
More generally the same summability and the local support condition
of \cref{thm:p15-bernoulli-product-complete} imply \(\mu\ll\nu\).
The endpoint extension asserts only the licensed direction.

The product condition alone does not repair failed local support:
if \(\mu\) has a fair first coordinate, \(\nu\) fixes its first
coordinate at zero, and their remaining independent coordinates
are identically fair, then \(\prod_ma_m=1/\sqrt2>0\) but
\(\mu\not\ll\nu\). These two measures are also not singular.
\end{hbtheorem}

\begin{proofboundary}{def:p15-density-exercise-scope}{The retained source exercise routes}
These results supply all Exercises~4.3.4--4.3.10 with their required
arguments. The event-product proof uses Borel--Cantelli; the density
proof uses finite partitions and an elementary upcrossing argument;
the Bernoulli comparisons use the full affinity criterion and a
separate event-occurrence counterexample. They preserve these
distinct routes even where an earlier general theorem already gives
the corresponding conclusion.
\end{proofboundary}

